\documentclass[11pt]{article}
\pagestyle{plain}

\usepackage[margin=1in]{geometry}
\usepackage{amsmath,amsfonts,amssymb,amsthm}
\usepackage[colorlinks=true,linkcolor=blue,citecolor=green,urlcolor=magenta]{hyperref}
\usepackage[capitalize,nameinlink]{cleveref}

\theoremstyle{plain}
\newtheorem{theorem}{Theorem}[section]
\newtheorem{proposition}[theorem]{Proposition}
\newtheorem{lemma}[theorem]{Lemma}
\newtheorem{claim}[theorem]{Claim}

\theoremstyle{definition}
\newtheorem{definition}[theorem]{Definition}
\newtheorem{example}[theorem]{Example}

\theoremstyle{remark}
\newtheorem{remark}[theorem]{Remark}

\newcommand{\R}{\mathbb{R}}
\newcommand{\abs}[1]{\left\lvert #1 \right\rvert}
\newcommand{\ip}[2]{\left\langle #1,#2 \right\rangle}
\newcommand{\norm}[1]{\left\lVert #1 \right\rVert}
\newcommand{\Hess}{\nabla^2}
\newcommand{\HessG}{\operatorname{Hess}^{\Gamma}}
\newcommand{\Ric}{\operatorname{Ric}}

\title{Curvature-Dimension Versus Reinforced Curvature-Dimension}
\date{\today}

\begin{document}
\maketitle

\section{Setup}

Let $L$ be the generator of a Markov semigroup. Recall the carré du champ
\begin{align*}
    \Gamma(f,g) := \frac{1}{2}\left(L(fg)-fLg-gLf\right),
    \qquad
    \Gamma(f) := \Gamma(f,f),
\end{align*}
and
\begin{align*}
    \Gamma_2(f,g) := \frac{1}{2}\left(L\Gamma(f,g)-\Gamma(f,Lg)-\Gamma(g,Lf)\right),
    \qquad
    \Gamma_2(f):=\Gamma_2(f,f).
\end{align*}
The curvature-dimension condition $\operatorname{CD}(\rho,\infty)$ is the pointwise inequality
\begin{align*}
    \Gamma_2(f) \ge \rho \Gamma(f)
\end{align*}
for all sufficiently regular $f$. The reinforced curvature-dimension condition used below is
\begin{align}\label{eq:rcd}
    \Gamma(\Gamma(f)) \le 4\Gamma(f)\left(\Gamma_2(f)-\rho\Gamma(f)\right).
\end{align}
The point of this note is that, for ordinary diffusion generators, $\operatorname{CD}(\rho,\infty)$ implies~\eqref{eq:rcd}; for discrete chains this implication is false.

\section{Diffusion Case}

In this section we assume that $L$ is a diffusion generator. Concretely, this
means that $\Gamma$ is a derivation in each argument:
\begin{align}\label{eq:derivation}
    \Gamma(ab,c)=a\Gamma(b,c)+b\Gamma(a,c).
\end{align}
Equivalently, for smooth functions one may think of
\[
    L=\sum_{i,j}a^{ij}\partial_{ij}+\sum_i b^i\partial_i,
    \qquad
    \Gamma(f,g)=\sum_{i,j}a^{ij}\partial_i f\,\partial_j g.
\]
Thus $\Gamma$ behaves like an inner product of gradients.

\begin{proposition}
Let $L$ be a diffusion generator satisfying $\operatorname{CD}(\rho,\infty)$.
Then the reinforced curvature-dimension inequality
\[
    \Gamma(\Gamma(f))
    \le
    4\Gamma(f)\left(\Gamma_2(f)-\rho\Gamma(f)\right)
\]
holds pointwise for every sufficiently regular $f$.
\end{proposition}

\begin{proof}
Fix a point $x$. All quantities below are evaluated at $x$. Define the
symmetric bilinear form
\[
    B(h,k):=\Gamma_2(h,k)-\rho\Gamma(h,k).
\]
The condition $\operatorname{CD}(\rho,\infty)$ says precisely that
\[
    B(h,h)\ge 0
\]
for every test function $h$. Hence, by Cauchy--Schwarz for positive
semidefinite bilinear forms,
\begin{align}\label{eq:B-cs}
    B(h,k)^2\le B(h,h)B(k,k).
\end{align}

Let $f$ and $g$ be test functions. Set
\[
    u:=f-f(x),\qquad v:=g-g(x),\qquad \eta:=uv.
\]
Then $u(x)=v(x)=0$. Since $\Gamma$ is a derivation,
\[
    \Gamma(\eta)(x)=0,
    \qquad
    \Gamma(f,\eta)(x)=0.
\]
Therefore
\[
    B(\eta,\eta)(x)=\Gamma_2(\eta)(x),
    \qquad
    B(f,\eta)(x)=\Gamma_2(f,\eta)(x).
\]
Applying \eqref{eq:B-cs} with $h=f$ and $k=\eta$ gives
\begin{align}\label{eq:first-cs}
    \Gamma_2(f,\eta)^2
    \le
    \left(\Gamma_2(f)-\rho\Gamma(f)\right)\Gamma_2(\eta).
\end{align}

We now compute the two terms involving $\eta$. First, using the derivation rule,
\[
    \Gamma(f,\eta)
    =
    \Gamma(f,uv)
    =
    u\Gamma(f,v)+v\Gamma(f,u).
\]
A standard expansion of $\Gamma_2$ using \eqref{eq:derivation} gives, at the
point $x$,
\begin{align}\label{eq:mixed-gamma2}
    \Gamma_2(f,\eta)
    =
    \Gamma(g,\Gamma(f)).
\end{align}
Similarly,
\begin{align}\label{eq:eta-gamma2}
    \Gamma_2(\eta)
    =
    2\left(\Gamma(f)\Gamma(g)+\Gamma(f,g)^2\right).
\end{align}
By Cauchy--Schwarz for the positive semidefinite form $\Gamma$,
\[
    \Gamma(f,g)^2\le \Gamma(f)\Gamma(g).
\]
Hence \eqref{eq:eta-gamma2} implies
\[
    \Gamma_2(\eta)
    \le
    4\Gamma(f)\Gamma(g).
\]
Combining this with \eqref{eq:first-cs} and \eqref{eq:mixed-gamma2}, we obtain
\begin{align}\label{eq:mixed-estimate}
    \Gamma(g,\Gamma(f))^2
    \le
    4\Gamma(f)\Gamma(g)
    \left(\Gamma_2(f)-\rho\Gamma(f)\right).
\end{align}

Finally choose $g=\Gamma(f)$. Then
\[
    \Gamma(g,\Gamma(f))=\Gamma(\Gamma(f)),
    \qquad
    \Gamma(g)=\Gamma(\Gamma(f)).
\]
Thus \eqref{eq:mixed-estimate} becomes
\[
    \Gamma(\Gamma(f))^2
    \le
    4\Gamma(f)\Gamma(\Gamma(f))
    \left(\Gamma_2(f)-\rho\Gamma(f)\right).
\]
If $\Gamma(\Gamma(f))(x)=0$, the desired inequality is immediate. Otherwise
we divide by $\Gamma(\Gamma(f))(x)$ and get
\[
    \Gamma(\Gamma(f))(x)
    \le
    4\Gamma(f)(x)
    \left(\Gamma_2(f)(x)-\rho\Gamma(f)(x)\right).
\]
Since $x$ was arbitrary, the reinforced inequality holds pointwise.
\end{proof}

\begin{remark}
The proof uses only the locality of a diffusion generator, encoded in the
derivation rule \eqref{eq:derivation}. This is precisely the step that fails
for jump generators on graphs. In the discrete case, even if
$\eta(x)=0$, the quantity $\Gamma(\eta)(x)$ may be positive because it depends
on values of $\eta$ at neighboring vertices.
\end{remark}

\section{A Discrete Counterexample}

Let $X=\{0,1,2\}$ and let $L$ be the continuous-time simple random walk generator on the path $0-1-2$:
\begin{align*}
    Lf(0)=f(1)-f(0),\qquad
    Lf(1)=f(0)+f(2)-2f(1),\qquad
    Lf(2)=f(1)-f(2).
\end{align*}
This chain is reversible with respect to the uniform measure. We use the same definitions of $\Gamma$ and $\Gamma_2$ as above.

\begin{lemma}
The path chain satisfies $\operatorname{CD}(1/2,\infty)$.
\end{lemma}

\begin{proof}
Write
\begin{align*}
    u=f(0)-f(1),\qquad v=f(2)-f(1).
\end{align*}
Adding a constant to $f$ does not change $\Gamma$ or $\Gamma_2$, so these two variables determine all the relevant quantities. Direct calculation gives
\begin{align*}
    \Gamma(f)(0)&=\frac{u^2}{2},
    &
    \Gamma_2(f)(0)&=u^2+\frac{uv}{2}+\frac{v^2}{4},
    \\
    \Gamma(f)(1)&=\frac{u^2+v^2}{2},
    &
    \Gamma_2(f)(1)&=\frac{3}{4}(u^2+v^2)+uv,
    \\
    \Gamma(f)(2)&=\frac{v^2}{2},
    &
    \Gamma_2(f)(2)&=v^2+\frac{uv}{2}+\frac{u^2}{4}.
\end{align*}
Therefore,
\begin{align*}
    \Gamma_2(f)(0)-\frac{1}{2}\Gamma(f)(0)
    &=
    \frac{(u+v)^2+2u^2}{4}
    \ge 0,\\
    \Gamma_2(f)(1)-\frac{1}{2}\Gamma(f)(1)
    &=
    \frac{(u+v)^2}{2}
    \ge 0,\\
    \Gamma_2(f)(2)-\frac{1}{2}\Gamma(f)(2)
    &=
    \frac{(u+v)^2+2v^2}{4}
    \ge 0.
\end{align*}
Thus $\Gamma_2(f)\ge \frac12 \Gamma(f)$ at every vertex for every $f$.
\end{proof}

\begin{proposition}
The same chain does not satisfy the reinforced inequality~\eqref{eq:rcd} with $\rho=1/2$.
\end{proposition}

\begin{proof}
Take
\begin{align*}
    f(0)=0,\qquad f(1)=0,\qquad f(2)=1.
\end{align*}
At the vertex $0$,
\begin{align*}
    \Gamma(f)(0)
    =
    \frac{1}{2}\left(f(1)-f(0)\right)^2
    =
    0.
\end{align*}
Also,
\begin{align*}
    \Gamma(f)(1)
    =
    \frac{1}{2}\left((f(0)-f(1))^2+(f(2)-f(1))^2\right)
    =
    \frac{1}{2}.
\end{align*}
Hence
\begin{align*}
    \Gamma(\Gamma(f))(0)
    =
    \frac{1}{2}\left(\Gamma(f)(1)-\Gamma(f)(0)\right)^2
    =
    \frac{1}{8}.
\end{align*}
The left-hand side of~\eqref{eq:rcd} at the vertex $0$ is therefore
\begin{align*}
    \Gamma(\Gamma(f))(0)=\frac18,
\end{align*}
whereas the right-hand side is
\begin{align*}
    4\Gamma(f)(0)\left(\Gamma_2(f)(0)-\frac12\Gamma(f)(0)\right)=0.
\end{align*}
Thus~\eqref{eq:rcd} fails.
\end{proof}

\begin{remark}
The counterexample is not caused by failure of the usual curvature-dimension condition: the previous lemma proves $\operatorname{CD}(1/2,\infty)$. The obstruction is that $\Gamma$ is nonlocal on a graph. Even if $f$ is locally constant across the only edge incident to $0$, so that $\Gamma(f)(0)=0$, the energy $\Gamma(f)$ can change across that edge because of edges adjacent to the neighboring vertex. Consequently $\Gamma(\Gamma(f))(0)$ can be positive.
\end{remark}

\end{document}
